\clearpage

% 6. RECURSOS
\section{Recursos Necessários}

\subsection{Membros do Grupo}

O grupo busca ter aproximadamente 8 a 15 membros ativos dispostos a:

\begin{itemize}
  \item Comparecer às sessões com regularidade;
  \item Contribuir para o compartilhamento de conhecimento e documentação;
  \item Participar de exercícios práticos;
  \item Colaborar construtivamente com outros membros.
\end{itemize}

\subsection{Espaço e Ambiente}

\begin{itemize}
  \item Computadores ou quadros para demonstrações e resolução colaborativa;
  \item Conexão estável à internet;
  \item Capacidade de projeção ou compartilhamento de tela para explicações.
\end{itemize}

\subsection{Plataformas de Aprendizado}

\begin{itemize}
  \item \textbf{Khan Academy} --- fundamentos e exercícios progressivos;
  \item \textbf{Brilliant.org} --- problemas interativos de matemática e lógica;
  \item \textbf{Art of Problem Solving (AoPS)} --- problemas voltados ao nível de competição;
  \item \textbf{CSES Problem Set} --- algoritmos com foco em programação;
  \item \textbf{Project Euler} --- problemas matemáticos resolvidos com programação.
\end{itemize}

\subsection{Ferramentas}

\begin{itemize}
  \item \textbf{Latex} --- um sistema de designer de documentos;
  \item \textbf{Markdown} --- uma linguagem de marcação de documentos HTML;
  \item \textbf{Git} --- um sistema de repositório de códigos (Como latex e markdown);
\end{itemize}
